\section{モデリングの目から見た検定5；カテゴリカル分布をつかって}
\subsection{授業内容}
\subsubsection{科目の中でのこのコマの位置づけ}
\subsubsection{コマ主題細目}
これまでは連続変数についてのモデルばかり扱ってきたが，度数など心理学で扱うデータの中にはカテゴリカルなものも少なくない。これらを帰無仮説検定の文脈で扱う時は，$\chi^2$ 検定が有用であるが，カテゴリカルな分布を使ったベイジアンアプローチももちろん，結果の解釈には有用かつ直感的である。

ここでは連続分布と離散分布の違いを確認し，いくつかの代表的な離散分布を演習とともに学び，カテゴリカルな出力変数の分析にも確率モデルが有用であることを確認する。
\begin{description}
	\item[離散的な分布] 変数には離散・連続の違いがあり，確率分布でも確率質量と確率密度の違いがある。ここでは代表的な離散確率変数であるベルヌーイ分布，二項分布，多項分布を導入する。

	      \begin{flushright}
		      $\to$\textcite{Matsuura201610}のPp.82，83，85--89.
	      \end{flushright}


	\item[$\chi^2$検定] カテゴリカルな変数の検定についてはこれまで扱って来なかったため，ここで改めて$\chi^2$分布を使った検定の例を導入する。$\chi^2$検定は比率(割合)，独立性，関連の強さ，適合度の検定などに用いられることを確認する。帰無仮説のおき方に注意することと，この検定がモデル適合度などの文脈でも望一られることに言及する。

	      \begin{flushright}
		      $\to$\textcite{Yamauchi201001}のPp.189-197.
	      \end{flushright}


	\item[カテゴリカル分布のモデリング] 確率分布を用いたアプローチをすることで，母比率を直接検証したり，連言命題が成立する確率など生成量を使ってさまざまな検証ができることを確認する。

	      \begin{flushright}
		      $\to$\textcite{Toyoda201606}のPp.136--163.
	      \end{flushright}

	\item[$\kappa$係数の算出] 変数間の関連の強さを見る指標として$\kappa$係数がある。一致率の係数ともして知られており，記述統計的アプローチでも算出できるものではあるが，確率モデルで表現する場合は工夫が必要である。

	      \begin{flushright}
		      $\to$\textcite{Lee201709}のPp.56--59.
	      \end{flushright}

\end{description}
\subsubsection{キーワード}
\begin{itemize}
	\item 離散分布と連続分布
	\item ベルヌーイ分布
	\item 二項分布
	\item 多項分布
	\item クロス集計表
	\item $\kappa$係数
\end{itemize}
\subsection{授業情報}
\paragraph{コマの展開方法}
講義/遠隔/演習
\subsubsection{予習・復習課題}
\paragraph{予習}確率変数が連続的か，離散的かということがどういう違いであるのかについて，データ解析基礎の確率に関する資料などを参考に確認しておく。また尺度水準による数値データの分類についても見直しておくと良い。
\paragraph{復習}カテゴリカルな分類，集計に関しては身の回りに多くのデータ例がある。たとえば官公庁の統計資料などをもとに，母比率の推定や連言命題が成立する確率など，さまざまな仮説を自ら立てて検証すると良い。